\documentclass[11pt,aspectratio=43]{beamer}

%------------------------------------------------
% Theme
%------------------------------------------------
\usetheme{Madrid}
\usefonttheme{serif}

%------------------------------------------------
% Colors
%------------------------------------------------
\definecolor{AegisBlue}{RGB}{52,50,181}
\setbeamercolor{structure}{fg=AegisBlue}
\setbeamercolor{frametitle}{bg=AegisBlue,fg=white}
\setbeamercolor{title}{bg=AegisBlue,fg=white}

%================================================
\begin{document}

%------------------------------------------------
% Slide 1 – Title Slide
%------------------------------------------------
\begin{frame}
\centering

% Title Box
%------------------------------------------------
% Blue Title Box
%------------------------------------------------
\begin{beamercolorbox}[
    wd=1\textwidth,
    rounded=true,
    shadow=true,
    sep=9pt,
    center
]{title}

\Large\bfseries
Aegis: A Secure and Autonomous
Delay-Tolerant Networking Framework
for Resilient Communication

\end{beamercolorbox}


\vspace{0.5cm}

\small
\begin{tabular}{l l}
Abhinav K P & TVE23CS003 \\
Akash B Mohan & TVE23CS016   \\
Alen K Raju & TVE23CS024 \\
Janbaas Jamal K K & TVE23CS071 \\
\end{tabular}

\vspace{0.4cm}


\normalsize \textbf{September 15, 2026}

\vspace{0.3cm}
\normalsize
\textbf{Project Guide}

\vspace{0.1cm}

\small
Prof.\ Rameez Mohammed A\\
Assistant Professor\\
Department of Computer Science and Engineering\\
College of Engineering Trivandrum

\end{frame}

%------------------------------------------------
% Slide 2 – Contents
%------------------------------------------------
\begin{frame}{Contents}
\tableofcontents
\end{frame}

%================================================
\section{Introduction}
%================================================

%------------------------------------------------
% Slide 3 – Introduction & Domain of Interest
%------------------------------------------------
\begin{frame}{Introduction}

\textbf{Delay-Tolerant Networking (DTN)} enables reliable communication in environments with intermittent connectivity by using the \textbf{Store--Carry--Forward} communication model.

\vspace{0.3cm}

\textbf{Domain of Interest}

\begin{itemize}
    \item Delay-Tolerant Networking (DTN)
    \item Embedded Systems \& FreeRTOS
    \item TinyML \& Edge AI
    \item Cybersecurity
    \item Space and Disaster Communication
\end{itemize}

\end{frame}

%================================================
\section{Problem Context}
%================================================

%------------------------------------------------
% Slide 4 – Problem Context
%------------------------------------------------
\begin{frame}{Problem Context}

\textbf{Mission-critical networks require reliable communication despite unstable connectivity and dynamic network conditions.}

\vspace{0.3cm}

\textbf{Key Challenges}

\begin{itemize}
    \item Intermittent links and unpredictable node failures
    \item Static DTN routing with limited adaptability
    \item No autonomous self-healing or intelligent route selection
    \item Poor priority handling and buffer utilization
    \item Limited trust-aware security for forwarding nodes
\end{itemize}

\end{frame}

%================================================
\section{Literature Survey}
%================================================
%------------------------------------------------
% Slide 14 – Paper 10
%------------------------------------------------
\begin{frame}{Literature Survey}

\textbf{Embedded AI: Integrating TinyML in Resource-Constrained IoT Devices}

\textit{IEEE Embedded Systems Research (2025)}

\vspace{0.2cm}

Presents lightweight TinyML deployment using TensorFlow Lite Micro, pruning, and quantization for resource-constrained embedded platforms.

\vspace{0.2cm}

\textbf{Merits}
\begin{itemize}
    \item Suitable for ESP32-class hardware
    \item Efficient model optimization
\end{itemize}

\textbf{Demerits}
\begin{itemize}
    \item No DTN integration
    \item No autonomous self-healing communication
\end{itemize}

\end{frame}
%------------------------------------------------
% Slide 13 – Paper 9
%------------------------------------------------
\begin{frame}{Literature Survey}

\textbf{TinyML for Edge Networks: Challenges and Future Directions}

\textit{IEEE ICCECE (2025)}

\vspace{0.2cm}

Reviews model compression, quantization, TensorFlow Lite Micro, and memory-efficient deployment techniques for embedded AI.

\vspace{0.2cm}

\textbf{Merits}
\begin{itemize}
    \item Optimized embedded AI deployment
    \item Low-memory inference techniques
\end{itemize}

\textbf{Demerits}
\begin{itemize}
    \item No networking application
    \item Does not address DTN routing
\end{itemize}

\end{frame}
%------------------------------------------------
% Slide 12 – Paper 8
%------------------------------------------------
\begin{frame}{Literature Survey}

\textbf{Distributed TinyML on Resource-Constrained IoT Sensor Networks}

\textit{IEEE World Forum on IoT (2024)}

\vspace{0.2cm}

Demonstrates distributed TinyML inference across multiple microcontrollers for collaborative edge intelligence in IoT systems.

\vspace{0.2cm}

\textbf{Merits}
\begin{itemize}
    \item TinyML on low-power MCUs
    \item Energy-efficient inference
\end{itemize}

\textbf{Demerits}
\begin{itemize}
    \item Focused on sensing applications
    \item No DTN communication framework
\end{itemize}

\end{frame}
%------------------------------------------------
% Slide 11 – Paper 7
%------------------------------------------------
\begin{frame}{Literature Survey}

\textbf{Optimal Routing with Machine Learning Classification in DTNs}

\textit{IEEE ICICNIS Conference (2024)}

\vspace{0.2cm}

Applies Random Forest, SVM, and XGBoost models to predict optimal forwarding nodes using DTN network parameters.

\vspace{0.2cm}

\textbf{Merits}
\begin{itemize}
    \item AI-assisted adaptive routing
    \item Higher delivery performance
\end{itemize}

\textbf{Demerits}
\begin{itemize}
    \item Evaluated only in simulation
    \item No embedded implementation
\end{itemize}

\end{frame}
%------------------------------------------------
% Slide 10 – Paper 6
%------------------------------------------------
\begin{frame}{Literature Survey}

\textbf{Spray and Wait Routing: Controlled Bundle Replication}

\textit{Spyropoulos, Psounis \& Raghavendra (IEEE, 2023)}

\vspace{0.2cm}

Reduces network congestion by limiting the number of replicated bundles before waiting for successful delivery.

\vspace{0.2cm}

\textbf{Merits}
\begin{itemize}
    \item Low bandwidth consumption
    \item Simple routing strategy
\end{itemize}

\textbf{Demerits}
\begin{itemize}
    \item Static forwarding policy
    \item No adaptive route selection
\end{itemize}

\end{frame}

%------------------------------------------------
% Slide 9 – Paper 5
%------------------------------------------------
\begin{frame}{Literature Survey}

\textbf{MaxProp Routing: Priority-Based Bundle Forwarding}

\textit{John Burgess et al. (IEEE INFOCOM, 2023)}

\vspace{0.2cm}

Prioritizes DTN bundles using delivery likelihood and acknowledgement history while managing limited buffer resources.

\vspace{0.2cm}

\textbf{Merits}
\begin{itemize}
    \item Priority-aware forwarding
    \item Improved delivery ratio
\end{itemize}

\textbf{Demerits}
\begin{itemize}
    \item High computational overhead
    \item Large buffer requirement
\end{itemize}

\end{frame}
%------------------------------------------------
% Slide 8 – Paper 4
%------------------------------------------------
\begin{frame}{Literature Survey}

\textbf{PRoPHET Routing: Probabilistic Routing in Intermittent Networks}

\textit{Anders Lindgren et al. (IEEE DTN Routing, 2023)}

\vspace{0.2cm}

Predicts successful bundle delivery using historical encounters between neighbouring DTN nodes to improve forwarding efficiency.

\vspace{0.2cm}

\textbf{Merits}
\begin{itemize}
    \item Learns encounter history
    \item Better than flooding-based routing
\end{itemize}

\textbf{Demerits}
\begin{itemize}
    \item Cannot predict sudden failures
    \item Ignores battery and buffer conditions
\end{itemize}

\end{frame}
%------------------------------------------------
% Slide 7 – Paper 3
%------------------------------------------------
\begin{frame}{Literature Survey}

\textbf{Contact Graph Routing (CGR): Predictive Routing for Space DTNs}

\textit{CCSDS / NASA Deep Space Networking (2022)}

\vspace{0.2cm}

Computes optimal forwarding paths using scheduled future communication opportunities between satellites and spacecraft.

\vspace{0.2cm}

\textbf{Merits}
\begin{itemize}
    \item Efficient deep-space routing
    \item Low routing overhead
\end{itemize}

\textbf{Demerits}
\begin{itemize}
    \item Requires predefined contact schedules
    \item Unsuitable for dynamic terrestrial DTNs
\end{itemize}

\end{frame}
%------------------------------------------------
% Slide 6 – Paper 2
%------------------------------------------------
\begin{frame}{Literature Survey}

\textbf{Bundle Protocol Security (BPsec): Security Framework for DTNs}

\textit{Edward Birrane, Kevin McKeever (RFC 9172, 2022)}

\vspace{0.2cm}

Introduces authentication, integrity verification, and confidential transmission for DTN bundles through cryptographic security blocks.

\vspace{0.2cm}

\textbf{Merits}
\begin{itemize}
    \item End-to-end authentication
    \item Secure bundle integrity
\end{itemize}

\textbf{Demerits}
\begin{itemize}
    \item No trust-aware routing
    \item Static security policies
\end{itemize}

\end{frame}

%------------------------------------------------
% Slide 5 – Paper 1
%------------------------------------------------
\begin{frame}{Literature Survey}

\textbf{Bundle Protocol Version 7 (BPv7): The Standard Protocol for Delay-Tolerant Networking}

\textit{Scott Burleigh, Kevin Fall, Edward Birrane (RFC 9171, 2022)}

\vspace{0.2cm}

Defines the standardized DTN communication protocol using the \textbf{Store--Carry--Forward} mechanism for reliable communication in disrupted networks.

\vspace{0.2cm}

\textbf{Merits}
\begin{itemize}
    \item Reliable bundle-based communication
    \item Supports intermittent connectivity
\end{itemize}

\textbf{Demerits}
\begin{itemize}
    \item No intelligent routing
    \item No autonomous self-healing
\end{itemize}

\end{frame}

%================================================
\section{Research Gap}
%================================================

%------------------------------------------------
% Slide 15 – Research Gap
%------------------------------------------------
\begin{frame}{Research Gap}

\textbf{Limitations Identified from Existing Literature}

\vspace{0.2cm}

\begin{itemize}
    \item DTN protocols rely on static or probabilistic routing.
    \item AI-based routing is mostly limited to simulations.
    \item TinyML research focuses on sensing rather than networking.
    \item Existing systems lack autonomous self-healing.
    \item No unified embedded DTN framework with security and AI.
\end{itemize}

\vspace{0.3cm}

\textbf{Our Contribution:} \textit{Aegis integrates DTN, TinyML, FreeRTOS, self-healing, priority-aware routing, and AES-256 security into a single ESP32-based embedded framework.}

\end{frame}
%================================================
\section{Problem Statement}
%================================================

%------------------------------------------------
% Slide 16 – Problem Statement
%------------------------------------------------
\begin{frame}{Problem Statement}

\vfill

\begin{block}{Problem}
\centering
The problem is to design and develop a secure, intelligent, and autonomous \textbf{Delay-Tolerant Networking (DTN)} framework that leverages \textbf{TinyML-assisted adaptive routing}, \textbf{self-healing communication}, and \textbf{priority-aware bundle management} to ensure reliable data delivery in environments with intermittent connectivity and dynamic node failures.
\end{block}

\vfill

\end{frame}
%================================================
\section{Project Objectives}
%================================================

%------------------------------------------------
% Slide 17 – Project Objectives
%------------------------------------------------
\begin{frame}{Project Objectives}

\begin{block}{Primary Objective}
Develop an intelligent, secure, and self-healing DTN framework on ESP32 using FreeRTOS and TinyML.
\end{block}

\vspace{0.2cm}

\textbf{Specific Objectives}
\begin{itemize}
    \item Implement AI-assisted adaptive routing
    \item Enable autonomous self-healing communication
    \item Prioritize mission-critical data bundles
    \item Optimize buffer and resource utilization
    \item Integrate AES-256 security and trust-based routing
\end{itemize}

\end{frame}

%================================================
\section{Proposed Solution Methodology}
%================================================

%------------------------------------------------
% Slide 17 – Overall Methodology
%------------------------------------------------
\begin{frame}{Proposed Solution Methodology}

\begin{block}{Workflow}
\small
Node Discovery $\rightarrow$ Network Monitoring $\rightarrow$ TinyML Prediction $\rightarrow$ Adaptive Route Selection $\rightarrow$ Priority Scheduling $\rightarrow$ Secure DTN Transmission
\end{block}

\vspace{0.3cm}

\textbf{Methodology Steps}
\begin{itemize}
    \item Discover neighbouring DTN nodes using heartbeat packets
    \item Continuously monitor link quality and node status
    \item Predict the best forwarding node using TinyML
    \item Prioritize mission-critical bundles before transmission
    \item Encrypt and forward bundles using Store--Carry--Forward
\end{itemize}

\end{frame}

%------------------------------------------------
% Slide 18 – Algorithm 1
%------------------------------------------------
\begin{frame}{Algorithm 1: Neighbor Discovery}

\textbf{Objective:} Identify and maintain active neighbouring DTN nodes.
\vspace{0.15cm}

\begin{figure}
    \centering
    \includegraphics[width=0.99\linewidth,height=5.9cm,keepaspectratio]{Neighbour Discovery.png}
    \caption{Neighbor Discovery Flow Diagram}
\end{figure}

\end{frame}


%------------------------------------------------
% Slide 19 – Algorithm 2
%------------------------------------------------
\begin{frame}{Algorithm 2: TinyML Adaptive Routing}

\textbf{Objective:} Predict the optimal next-hop node.

\vspace{0.2cm}

\textbf{Inputs:} RSSI, Latency, Buffer, Battery, Trust

\vspace{0.15cm}

\begin{figure}
    \centering
    \includegraphics[width=0.72\linewidth,height=5.0cm,keepaspectratio]{TinyML.png}
    \caption{TinyML-Based Adaptive Routing Flowchart}
\end{figure}

\end{frame}


%------------------------------------------------
% Slide 20 – Algorithm 3
%------------------------------------------------
\begin{frame}{Algorithm 3: Self-Healing Mechanism}

\textbf{Objective:} Restore communication after node failure.
\vspace{0.15cm}

\begin{figure}
    \centering
    \includegraphics[width=0.99\linewidth,height=5.9cm,keepaspectratio]{Self Healing.png}
    \caption{Self-Healing Communication Mechanism}
\end{figure}

\end{frame}
%================================================
\section{System Architecture}
%================================================

%------------------------------------------------
% Slide 21 – System Architecture
%------------------------------------------------
\begin{frame}{System Architecture}

\begin{figure}
    \centering
    \includegraphics[height=0.78\textheight,keepaspectratio]{Architecturev2.png}
    \caption{Proposed System Architecture}
\end{figure}

\end{frame}

%------------------------------------------------
% Slide 22 – Block Diagram
%------------------------------------------------
\begin{frame}{Block Diagram of Proposed System}

\begin{figure}
    \centering
    \includegraphics[width=1\linewidth,height=7.3cm,keepaspectratio]{block.png}
    \caption{Block Diagram of the Proposed DTN System}
\end{figure}

\end{frame}
%================================================
\section{Project Planning}
%================================================

%------------------------------------------------
% Slide 23 – Project Planning
%------------------------------------------------
\begin{frame}[t]{Project Planning \& Task Allocation}

\small
\setlength{\tabcolsep}{5pt}
\renewcommand{\arraystretch}{1.25}

\begin{block}{1. Project Schedule}
\scriptsize
\begin{tabular*}{\linewidth}{@{\extracolsep{\fill}}llccccccc}

\textcolor{black}{\textbf{Phase}} &
\textcolor{black}{\textbf{Activity}} &
\textcolor{black}{\textbf{Sep}} &
\textcolor{black}{\textbf{Oct}} &
\textcolor{black}{\textbf{Nov}} &
\textcolor{black}{\textbf{Dec}} &
\textcolor{black}{\textbf{Jan}} &
\textcolor{black}{\textbf{Feb}} &
\textcolor{black}{\textbf{Mar}} \\

\textbf{Planning} & Literature Survey
    & \cellcolor{-} & & & & & & \\
& System Design
    & \cellcolor{-} & \cellcolor{-} & & & & & \\ \hline

\textbf{Development} & DTN Protocol Implementation
    & & \cellcolor{-} & \cellcolor{-} & & & & \\
& TinyML Model Development
    & & & \cellcolor{-} & \cellcolor{-} & & & \\
& ESP32 + FreeRTOS Integration
    & & & & \cellcolor{-} & \cellcolor{-} & & \\
& Self-Healing Routing
    & & & & & \cellcolor{-} & \cellcolor{-} & \\ \hline

\textbf{Testing} & System Integration \& Testing
    & & & & & & \cellcolor{-} & \cellcolor{-} \\
& Performance Evaluation
    & & & & & & \cellcolor{-} & \cellcolor{-} \\

\end{tabular*}
\end{block}

\begin{block}{2. Team Responsibilities}
\scriptsize
\begin{tabular*}{\linewidth}{@{\extracolsep{\fill}}ll}

\textcolor{black}{\textbf{Team Member}} &
\textcolor{black}{\textbf{Primary Responsibility}} \\


\textbf{Abhinav K P} & TinyML \& System Architecture \\
\textbf{Akash B Mohan} & Self-Healing Routing \& Literature Survey \\
\textbf{Alen K Raju} & Security Layer \& System Integration \\
\textbf{Janbaas Jamal K K} & ESP32, FreeRTOS \& DTN Communication \\
\end{tabular*}
\end{block}

\end{frame}

%================================================
\section{Conclusion}
%================================================

%------------------------------------------------
% Slide 24 – Conclusion
%------------------------------------------------
\begin{frame}{Conclusion}

\begin{itemize}
    \item Proposed an intelligent and secure embedded DTN framework for disruption-prone networks.
    \item Integrates TinyML-based adaptive routing with autonomous self-healing communication.
    \item Enables priority-aware bundle management and AES-256 secure data transmission.
    \item Targets real-time deployment on ESP32 using FreeRTOS for resilient IoT applications.
\end{itemize}

\end{frame}
%================================================
\section{References}
%================================================

%------------------------------------------------
% Slide 24 – References
%------------------------------------------------
\begin{frame}{References}
\tiny
\setlength{\baselineskip}{8.8pt}

[1] ``TinyML for Edge Networks: Challenges and Future Directions,'' V. M. R. Tummala, S. Korada, S. P. Lingamallu, B. S. Hari, and A. Hazra, \textit{Proc. IEEE ICCECE}, 2025.\\

[2] ``Embedded AI: Integrating TinyML in Resource-Constrained IoT Devices,'' M. Sharma, R. Patel, and K. Verma, \textit{Proc. IEEE IMC-SSGP}, 2025.\\

[3] ``TinyML on Microcontrollers: Energy-Efficient Real-Time Incremental Learning,'' A. Kumar, P. Nair, and S. Reddy, \textit{Proc. IEEE ICSCN}, 2025.\\

[4] ``Distributed TinyML on Resource-Constrained IoT Sensor Networks,'' J. Smith, L. Chen, and R. Brown, \textit{Proc. IEEE WF-IoT}, 2024.\\

[5] ``Optimal Routing with Machine Learning Classification in Delay-Tolerant Networks,'' Y. Zhang, H. Liu, and X. Wang, \textit{Proc. IEEE ICICNIS}, 2024.\\

[6] ``PRoPHET Routing Protocol for Intermittently Connected Networks,'' A. Lindgren, A. Doria, and O. Schel\'en, \textit{IEEE Trans. Mobile Computing}, 2023.\\

[7] ``MaxProp: Routing for Vehicle-Based Delay-Tolerant Networks,'' J. Burgess, B. Gallagher, D. Jensen, and B. N. Levine, \textit{Proc. IEEE INFOCOM}, 2023.\\

[8] ``Spray and Wait: An Efficient Routing Scheme for Delay-Tolerant Networks,'' T. Spyropoulos, K. Psounis, and C. S. Raghavendra, \textit{Proc. IEEE}, 2023.\\

[9] ``Bundle Protocol Security (BPsec),'' E. Birrane and K. McKeever, \textit{IETF RFC 9172}, 2022.\\

[10] ``Bundle Protocol Version 7,'' S. Burleigh, K. Fall, and E. Birrane, \textit{IETF RFC 9171}, 2022.

\end{frame}
\end{document}